\section{Двойственность и корреляция}

\begin{definition}
    Пусть $v \in V,  \text{определим} \langle -, v\rangle: V \rightarrow \mathbb{R}$.
    Эта функция линейна по первому аргкменту $\Rightarrow \langle -, v\rangle \in V^*$
    \[
    G_V \colon V \to V^*, \quad v \mapsto \langle-, v\rangle
    \]
\end{definition}

\begin{remark}
    Не путать с матрицей Грама, тут $V$ -- большое.
\end{remark}

\begin{stmt}
    $G_V$ --- изоморфизм (если $V$ конечномерно, как и в многих утверждениях раньше).
    \begin{proof}
        \textbf{Инъективность.} $\ker G_V = \{v \mid \langle -, v \rangle \equiv 0\} = \{0\}$,
        поскольку при $v \neq 0$ функционал $\langle -, v \rangle$ ненулевой:
        $\langle v, v \rangle > 0$.

        \textbf{Сюръективность.} $\dim V = \dim V^*$, поэтому инъекция автоматически является биекцией.

        \textbf{Линейность.} Скалярное произведение линейно по второму аргументу,
        откуда $G_V$ линейно. \\ Следовательно, $G_V$ --- изоморфизм.
    \end{proof}

\end{stmt}

\begin{remark}
    Такой изоморфизм \textit{почти канонический}.
\end{remark}

\begin{definition}[$G_V$ — евклидова корреляция]
    Пусть $e_1, \dots, e_n$ — базис $V$ и $e_1^*, \dots, e_n^* \in V^*$ — двойств.\ базис.
    Положим
    \[
        e_i^\times := G_V^{-1}(e_i^*) \in V, \qquad
        \langle e_i, e_j^\times \rangle = e_i^*(e_j) = \begin{cases} 1, & i = j, \\ 0, & i \neq j. \end{cases}
    \]

\end{definition}

\begin{definition}
    $e_1^\times, \dots, e_n^\times$ называется \emph{евклидовым двойственным базисом} пространства $V$.

    $G_{\mathbf{e} \mathbf{e^\times}} = E$ — матрица Грама наборов $\mathbf{e}$ и $\mathbf{e^\times}$.
\end{definition}

\begin{remark}
    Пусть $C_{\mathbf{e}\,\mathbf{e}^\times}$~--- матрица перехода от $\mathbf{e}$ к
    $\mathbf{e}^\times$. Тогда
    \[
        G_{\mathbf{e}\,\mathbf{e}^\times} = G_{\mathbf{e}} \cdot C_{\mathbf{e}\,\mathbf{e}^\times}.
    \]
\end{remark}

\begin{proplist}
    \item
    Если $e$ — ортонорм.\ базис, то $\mathbf{e^\times} = \mathbf{e}$.
    \item
    $(\mathbf{e^\times})^\times = \mathbf{e}$
\begin{proof}
    \begin{proofsteps}
        \item Покажем, что $G_V(e_i) = e_i^*$, тогда $e_i = G_V^{-1}(e_i^*) = e_i^\times$:
        \[
            (G_V e_i)(e_j)
            = \langle e_j, e_i \rangle
            = \delta_{ij}
            = e_i^*(e_j).
        \]
        Следовательно, $e_i^\times = e_i$ для всех $i$.

        \item Покажем, что $e_i = e_i^{\times\times}$, т.\,е.\ $G_V(e_i) = (e_i^\times)^*$:
        \[
            G_V(e_i)(e_j^\times)
            = \langle e_j^\times, e_i \rangle
            = \langle e_i, e_j^\times \rangle
            = \delta_{ij}
            = (e_i^\times)^*(e_j^\times).
        \]
        Где $\langle e_i, e_j^\times \rangle = \delta_{ij}$ — по определению $e_j^\times$.
        Следовательно, $e_i = G_V^{-1}((e_i^\times)^*) = e_i^{\times\times}$.
    \end{proofsteps}
\end{proof}

\end{proplist}
